\section{AHSME 1955}
        \begin{enumerate}
        	\item (\textit{AHSME 1955}) Which one of the following is not equivalent to $0.000000375$? 


 $\textbf{(A)}\ 3.75\times 10^{-7}\qquad\textbf{(B)}\ 3\frac{3}{4}\times 10^{-7}\qquad\textbf{(C)}\ 375\times 10^{-9}\\ \textbf{(D)}\ \frac{3}{8}\times 10^{-7}\qquad\textbf{(E)}\ \frac{3}{80000000}$


	\item (\textit{AHSME 1955}) The smaller angle between the hands of a clock at $12:25$ p.m. is: 


 $\textbf{(A)}\ 132^\circ 30'\qquad\textbf{(B)}\ 137^\circ 30'\qquad\textbf{(C)}\ 150^\circ\qquad\textbf{(D)}\ 137^\circ 32'\qquad\textbf{(E)}\ 137^\circ$


	\item (\textit{AHSME 1955}) If each number in a set of ten numbers is increased by $20$, the arithmetic mean (average) of the ten numbers: 


 $\textbf{(A)}\ \text{remains the same}\qquad\textbf{(B)}\ \text{is increased by 20}\qquad\textbf{(C)}\ \text{is increased by 200}\\ \textbf{(D)}\ \text{is increased by 10}\qquad\textbf{(E)}\ \text{is increased by 2}$


	\item (\textit{AHSME 1955}) The equality $\frac{1}{x-1}=\frac{2}{x-2}$ is satisfied by: 


 $\textbf{(A)}\ \text{no real values of }x\qquad\textbf{(B)}\ \text{either }x=1\text{ or }x=2\qquad\textbf{(C)}\ \text{only }x=1\\ \textbf{(D)}\ \text{only }x=2\qquad\textbf{(E)}\ \text{only }x=0$


	\item (\textit{AHSME 1955}) $5y$ varies inversely as the square of $x$. When $y=16, x=1$. When $x=8, y$ equals: 


 $\textbf{(A)}\ 2 \qquad \textbf{(B)}\ 128 \qquad \textbf{(C)}\ 64 \qquad \textbf{(D)}\ \frac{1}{4} \qquad \textbf{(E)}\ 1024$


	\item (\textit{AHSME 1955}) A merchant buys a number of oranges at $3$ for $10$ cents and an equal number at $5$ for $20$ cents. To "break even" he must sell all at: 


 $\textbf{(A)}\ \text{8 for 30 cents}\qquad\textbf{(B)}\ \text{3 for 11 cents}\qquad\textbf{(C)}\ \text{5 for 18 cents}\\ \textbf{(D)}\ \text{11 for 40 cents}\qquad\textbf{(E)}\ \text{13 for 50 cents}$


	\item (\textit{AHSME 1955}) If a worker receives a $20$\% cut in wages, he may regain his original pay exactly by obtaining a raise of:


 $\textbf{(A)}\ \text{20\%}\qquad\textbf{(B)}\ \text{25\%}\qquad\textbf{(C)}\ 22\frac{1}{2}\text{\%}\qquad\textbf{(D)}\ \textdollar{20}\qquad\textbf{(E)}\ \textdollar{25}$


	\item (\textit{AHSME 1955}) The graph of $x^2-4y^2=0$: 


 $\textbf{(A)}\ \text{is a hyperbola intersecting only the }x\text{-axis}\\ \textbf{(B)}\ \text{is a hyperbola intersecting only the }y\text{-axis}\\ \textbf{(C)}\ \text{is a hyperbola intersecting neither axis}\\ \textbf{(D)}\ \text{is a pair of straight lines}\\ \textbf{(E)}\ \text{does not exist}$


	\item (\textit{AHSME 1955}) A circle is inscribed in a triangle with sides $8, 15$, and $17$. The radius of the circle is: 


 $\textbf{(A)}\ 6 \qquad \textbf{(B)}\ 2 \qquad \textbf{(C)}\ 5 \qquad \textbf{(D)}\ 3 \qquad \textbf{(E)}\ 7$


	\item (\textit{AHSME 1955}) How many hours does it take a train traveling at an average rate of 40 mph between stops to travel a miles it makes n stops of m minutes each? 


 $\textbf{(A)}\ \frac{3a+2mn}{120}\qquad\textbf{(B)}\ 3a+2mn\qquad\textbf{(C)}\ \frac{3a+2mn}{12}\qquad\textbf{(D)}\ \frac{a+mn}{40}\qquad\textbf{(E)}\ \frac{a+40mn}{40}$


	\item (\textit{AHSME 1955}) The negation of the statement "No slow learners attend this school" is: 


 $\textbf{(A)}\ \text{All slow learners attend this school}\\ \textbf{(B)}\ \text{All slow learners do not attend this school}\\ \textbf{(C)}\ \text{Some slow learners attend this school}\\ \textbf{(D)}\ \text{Some slow learners do not attend this school}\\ \textbf{(E)}\ \text{No slow learners do not attend this school}$


	\item (\textit{AHSME 1955}) The solution of $\sqrt{5x-1}+\sqrt{x-1}=2$ is: 


 $\textbf{(A)}\ x=2,x=1\qquad\textbf{(B)}\ x=\frac{2}{3}\qquad\textbf{(C)}\ x=2\qquad\textbf{(D)}\ x=1\qquad\textbf{(E)}\ x=0$


	\item (\textit{AHSME 1955}) The fraction $\frac{a^{-4}-b^{-4}}{a^{-2}-b^{-2}}$ is equal to: 


 $\textbf{(A)}\ a^{-6}-b^{-6}\qquad\textbf{(B)}\ a^{-2}-b^{-2}\qquad\textbf{(C)}\ a^{-2}+b^{-2}\\ \textbf{(D)}\ a^2+b^2\qquad\textbf{(E)}\ a^2-b^2$


	\item (\textit{AHSME 1955}) The length of rectangle $R$ is $10$\% more than the side of square $S$. The width of the rectangle is $10$\% less than the side of the square. 
The ratio of the areas, $R:S$, is: 


 $\textbf{(A)}\ 99: 100\qquad\textbf{(B)}\ 101: 100\qquad\textbf{(C)}\ 1: 1\qquad\textbf{(D)}\ 199: 200\qquad\textbf{(E)}\ 201: 200$


	\item (\textit{AHSME 1955}) The ratio of the areas of two concentric circles is $1: 3$. If the radius of the smaller is $r$, then the difference between the 
radii is best approximated by: 


 $\textbf{(A)}\ 0.41r \qquad \textbf{(B)}\ 0.73 \qquad \textbf{(C)}\ 0.75 \qquad \textbf{(D)}\ 0.73r \qquad \textbf{(E)}\ 0.75r$


	\item (\textit{AHSME 1955}) The value of $\frac{3}{a+b}$ when $a=4$ and $b=-4$ is: 


 $\textbf{(A)}\ 3\qquad\textbf{(B)}\ \frac{3}{8}\qquad\textbf{(C)}\ 0\qquad\textbf{(D)}\ \text{any finite number}\qquad\textbf{(E)}\ \text{meaningless}$


	\item (\textit{AHSME 1955}) If $\log x-5 \log 3=-2$, then $x$ equals: 


 $\textbf{(A)}\ 1.25\qquad\textbf{(B)}\ 0.81\qquad\textbf{(C)}\ 2.43\qquad\textbf{(D)}\ 0.8\qquad\textbf{(E)}\ \text{either 0.8 or 1.25}$


	\item (\textit{AHSME 1955}) The discriminant of the equation $x^2+2x\sqrt{3}+3=0$ is zero. Hence, its roots are: 


 $\textbf{(A)}\ \text{real and equal}\qquad\textbf{(B)}\ \text{rational and equal}\qquad\textbf{(C)}\ \text{rational and unequal}\\ \textbf{(D)}\ \text{irrational and unequal}\qquad\textbf{(E)}\ \text{imaginary}$


	\item (\textit{AHSME 1955}) Two numbers whose sum is $6$ and the absolute value of whose difference is $8$ are roots of the equation: 


 $\textbf{(A)}\ x^2-6x+7=0\qquad\textbf{(B)}\ x^2-6x-7=0\qquad\textbf{(C)}\ x^2+6x-8=0\\ \textbf{(D)}\ x^2-6x+8=0\qquad\textbf{(E)}\ x^2+6x-7=0$


	\item (\textit{AHSME 1955}) The expression $\sqrt{25-t^2}+5$ equals zero for: 


 $\textbf{(A)}\ \text{no real or imaginary values of }t\qquad\textbf{(B)}\ \text{no real values of }t\text{ only}\\ \textbf{(C)}\ \text{no imaginary values of }t\text{ only}\qquad\textbf{(D)}\ t=0\qquad\textbf{(E)}\ t=\pm 5$


	\item (\textit{AHSME 1955}) Represent the hypotenuse of a right triangle by $c$ and the area by $A$. The altitude on the hypotenuse is: 


 $\textbf{(A)}\ \frac{A}{c}\qquad\textbf{(B)}\ \frac{2A}{c}\qquad\textbf{(C)}\ \frac{A}{2c}\qquad\textbf{(D)}\ \frac{A^2}{c}\qquad\textbf{(E)}\ \frac{A}{c^2}$


	\item (\textit{AHSME 1955}) On a $\textdollar{10000}$ order a merchant has a choice between three successive discounts of $20$\%, $20$\%, and $10$\% and 
three successive discounts of $40$\%, $5$\%, and $5$\%. By choosing the better offer, he can save: 


 $\textbf{(A)}\ \text{nothing at all}\qquad\textbf{(B)}\ \$440\qquad\textbf{(C)}\ \$330\qquad\textbf{(D)}\ \$345\qquad\textbf{(E)}\ \$360$


	\item (\textit{AHSME 1955}) In checking the petty cash a clerk counts $q$ quarters, $d$ dimes, $n$ nickels, and $c$ cents. 
Later he discovers that $x$ of the nickels were counted as quarters and $x$ of the dimes were counted as cents. 
To correct the total obtained the clerk must: 


 $\textbf{(A)}\ \text{make no correction}\qquad\textbf{(B)}\ \text{subtract 11 cents}\qquad\textbf{(C)}\ \text{subtract 11}x\text{ cents}\\ \textbf{(D)}\ \text{add 11}x\text{ cents}\qquad\textbf{(E)}\ \text{add }x\text{ cents}$


	\item (\textit{AHSME 1955}) The function $4x^2-12x-1$: 


 $\textbf{(A)}\ \text{always increases as }x\text{ increases}\\ \textbf{(B)}\ \text{always decreases as }x\text{ decreases to 1}\\ \textbf{(C)}\ \text{cannot equal 0}\\ \textbf{(D)}\ \text{has a maximum value when }x\text{ is negative}\\ \textbf{(E)}\ \text{has a minimum value of-10}$


	\item (\textit{AHSME 1955}) One of the factors of $x^4+2x^2+9$ is: 


 $\textbf{(A)}\ x^2+3\qquad\textbf{(B)}\ x+1\qquad\textbf{(C)}\ x^2-3\qquad\textbf{(D)}\ x^2-2x-3\qquad\textbf{(E)}\ \text{none of these}$


	\item (\textit{AHSME 1955}) Mr. A owns a house worth $\textdollar{10000}$. He sells it to Mr. $B$ at $10$\% profit. Mr. $B$ sells the house back to Mr. $A$ at a $10$\% loss. Then: 


 $\textbf{(A)}\ \text{Mr. A comes out even}\qquad\textbf{(B)}\ \text{Mr. A makes } \textdollar{ 100}\qquad\textbf{(C)}\ \text{Mr. A makes } \textdollar{ 1000}\\ \textbf{(D)}\ \text{Mr. B loses } \textdollar{ 100}\qquad\textbf{(E)}\ \text{none of the above is correct}$


	\item (\textit{AHSME 1955}) If $r$ and $s$ are the roots of $x^2-px+q=0$, then $r^2+s^2$ equals: 


 $\textbf{(A)}\ p^2+2q\qquad\textbf{(B)}\ p^2-2q\qquad\textbf{(C)}\ p^2+q^2\qquad\textbf{(D)}\ p^2-q^2\qquad\textbf{(E)}\ p^2$


	\item (\textit{AHSME 1955}) On the same set of axes are drawn the graph of $y=ax^2+bx+c$ and the graph of the equation obtained by replacing $x$ by $-x$ in the given equation. 
If $b \neq 0$ and $c \neq 0$ these two graphs intersect: 


 $\textbf{(A)}\ \text{in two points, one on the x-axis and one on the y-axis}\\ \textbf{(B)}\ \text{in one point located on neither axis}\\ \textbf{(C)}\ \text{only at the origin}\\ \textbf{(D)}\ \text{in one point on the x-axis}\\ \textbf{(E)}\ \text{in one point on the y-axis}$


	\item (\textit{AHSME 1955}) In the figure, $PA$ is tangent to semicircle $SAR$; $PB$ is tangent to semicircle $RBT$; $SRT$ is a straight line; 
the arcs are indicated in the figure. $\angle APB$ is measured by: 


 
                    \begin{figure}[H]
                        \centering
                        \includegraphics[width=0.5\textwidth]{Problem/AHSME_I_1955_29_Problem_Figure_1.png}
                    \end{figure}
                    
$\textbf{(A)}\ \frac{1}{2}(a-b)\qquad\textbf{(B)}\ \frac{1}{2}(a+b)\qquad\textbf{(C)}\ (c-a)-(d-b)\qquad\textbf{(D)}\ a-b\qquad\textbf{(E)}\ a+b$


	\item (\textit{AHSME 1955}) Each of the equations $3x^2-2=25, (2x-1)^2=(x-1)^2, \sqrt{x^2-7}=\sqrt{x-1}$ has: 


 $\textbf{(A)}\ \text{two integral roots}\qquad\textbf{(B)}\ \text{no root greater than 3}\qquad\textbf{(C)}\ \text{no root zero}\\ \textbf{(D)}\ \text{only one root}\qquad\textbf{(E)}\ \text{one negative root and one positive root}$


	\item (\textit{AHSME 1955}) An equilateral triangle whose side is $2$ is divided into a triangle and a trapezoid by a line drawn parallel to one of its sides. 
If the area of the trapezoid equals one-half of the area of the original triangle, the length of the median of the trapezoid is: 


 $\textbf{(A)}\ \frac{\sqrt{6}}{2}\qquad\textbf{(B)}\ \sqrt{2}\qquad\textbf{(C)}\ 2+\sqrt{2}\qquad\textbf{(D)}\ \frac{2+\sqrt{2}}{2}\qquad\textbf{(E)}\ \frac{2\sqrt{3}-\sqrt{6}}{2}$


	\item (\textit{AHSME 1955}) If the discriminant of $ax^2+2bx+c=0$ is zero, then another true statement about $a, b$, and $c$ is that: 


 $\textbf{(A)}\ \text{they form an arithmetic progression}\\ \textbf{(B)}\ \text{they form a geometric progression}\\ \textbf{(C)}\ \text{they are unequal}\\ \textbf{(D)}\ \text{they are all negative numbers}\\ \textbf{(E)}\ \text{only b is negative and a and c are positive}$


	\item (\textit{AHSME 1955}) Henry starts a trip when the hands of the clock are together between $8$ a.m. and $9$ a.m. 
He arrives at his destination between $2$ p.m. and $3$ p.m. when the hands of the clock are exactly $180^\circ$ apart. The trip takes: 


 $\textbf{(A)}\ \text{6 hr.}\qquad\textbf{(B)}\ \text{6 hr.   } 43\frac{7}{11} \text{min.}\qquad\textbf{(C)}\ \text{5 hr.   } 16\frac{4}{11} \text{min.}\qquad\textbf{(D)}\ \text{6 hr. 30 min.}\qquad\textbf{(E)}\ \text{none of these}$


	\item (\textit{AHSME 1955}) A $6$-inch and $18$-inch diameter pole are placed together and bound together with wire. 
The length of the shortest wire that will go around them is: 


 $\textbf{(A)}\ 12\sqrt{3}+16\pi\qquad\textbf{(B)}\ 12\sqrt{3}+7\pi\qquad\textbf{(C)}\ 12\sqrt{3}+14\pi\\ \textbf{(D)}\ 12+15\pi\qquad\textbf{(E)}\ 24\pi$


	\item (\textit{AHSME 1955}) Three boys agree to divide a bag of marbles in the following manner. The first boy takes one more than half the marbles. 
The second takes a third of the number remaining. The third boy finds that he is left with twice as many marbles as the second boy. 
The original number of marbles: 


 $\textbf{(A)}\ \text{is none of the following}\qquad\textbf{(B)}\ \text{cannot be determined from the given data}\\ \textbf{(C)}\ \text{is 20 or 26}\qquad\textbf{(D)}\ \text{is 14 or 32}\qquad\textbf{(E)}\ \text{is 8 or 38}$


	\item (\textit{AHSME 1955}) A cylindrical oil tank, lying horizontally, has an interior length of $10$ feet and an interior diameter of $6$ feet. 
If the rectangular surface of the oil has an area of $40$ square feet, the depth of the oil is: 


 $\textbf{(A)}\ \sqrt{5}\qquad\textbf{(B)}\ 2\sqrt{5}\qquad\textbf{(C)}\ 3-\sqrt{5}\qquad\textbf{(D)}\ 3+\sqrt{5}\\ \textbf{(E)}\ \text{either }3-\sqrt{5}\text{ or }3+\sqrt{5}$


	\item (\textit{AHSME 1955}) A three-digit number has, from left to right, the digits $h, t$, and $u$, with $h>u$. 
When the number with the digits reversed is subtracted from the original number, the units' digit in the difference is 4. 
The next two digits, from right to left, are: 


 $\textbf{(A)}\ \text{5 and 9}\qquad\textbf{(B)}\ \text{9 and 5}\qquad\textbf{(C)}\ \text{impossible to tell}\qquad\textbf{(D)}\ \text{5 and 4}\qquad\textbf{(E)}\ \text{4 and 5}$


	\item (\textit{AHSME 1955}) Four positive integers are given. Select any three of these integers, find their arithmetic average, 
and add this result to the fourth integer. Thus the numbers $29, 23, 21$, and $17$ are obtained. One of the original integers is: 


 $\textbf{(A)}\ 19 \qquad \textbf{(B)}\ 21 \qquad \textbf{(C)}\ 23 \qquad \textbf{(D)}\ 29 \qquad \textbf{(E)}\ 17$


	\item (\textit{AHSME 1955}) If $y=x^2+px+q$, then if the least possible value of $y$ is zero $q$ is equal to: 


 $\textbf{(A)}\ 0\qquad\textbf{(B)}\ \frac{p^2}{4}\qquad\textbf{(C)}\ \frac{p}{2}\qquad\textbf{(D)}\ -\frac{p}{2}\qquad\textbf{(E)}\ \frac{p^2}{4}-q$


	\item (\textit{AHSME 1955}) The fractions $\frac{ax+b}{cx+d}$ and $\frac{b}{d}$ are unequal if: 


 $\textbf{(A)}\ a=c=1, x\neq 0\qquad\textbf{(B)}\ a=b=0\qquad\textbf{(C)}\ a=c=0\\ \textbf{(D)}\ x=0\qquad\textbf{(E)}\ ad=bc$


	\item (\textit{AHSME 1955}) A train traveling from Aytown to Beetown meets with an accident after $1$ hr. It is stopped for $\frac{1}{2}$ hr., 
after which it proceeds at four-fifths of its usual rate, arriving at Beetown $2$ hr. late. 
If the train had covered $80$ miles more before the accident, it would have been just $1$ hr. late. 
The usual rate of the train is: 


 $\textbf{(A)}\ \text{20 mph}\qquad\textbf{(B)}\ \text{30 mph}\qquad\textbf{(C)}\ \text{40 mph}\qquad\textbf{(D)}\ \text{50 mph}\qquad\textbf{(E)}\ \text{60 mph}$


	\item (\textit{AHSME 1955}) If $a, b$, and $c$ are positive integers, the radicals $\sqrt{a+\frac{b}{c}}$ and $a\sqrt{\frac{b}{c}}$ are equal when and only when: 


 $\textbf{(A)}\ a=b=c=1\qquad\textbf{(B)}\ a=b\text{ and }c=a=1\qquad\textbf{(C)}\ c=\frac{b(a^2-1)}{a}\\ \textbf{(D)}\ a=b\text{ and }c\text{ is any value}\qquad\textbf{(E)}\ a=b\text{ and }c=a-1$


	\item (\textit{AHSME 1955}) The pairs of values of $x$ and $y$ that are the common solutions of the equations $y=(x+1)^2$ and $xy+y=1$ are: 


 $\textbf{(A)}\ \text{3 real pairs}\qquad\textbf{(B)}\ \text{4 real pairs}\qquad\textbf{(C)}\ \text{4 imaginary pairs}\\ \textbf{(D)}\ \text{2 real and 2 imaginary pairs}\qquad\textbf{(E)}\ \text{1 real and 2 imaginary pairs}$


	\item (\textit{AHSME 1955}) In circle $O$ chord $AB$ is produced so that $BC$ equals a radius of the circle. $CO$ is drawn and extended to $D$. 
$AO$ is drawn. Which of the following expresses the relationship between $x$ and $y$?


 
                    \begin{figure}[H]
                        \centering
                        \includegraphics[width=0.5\textwidth]{Problem/AHSME_I_1955_44_Problem_Figure_1.png}
                    \end{figure}
                    
$\textbf{(A)}\ x=3y\\ \textbf{(B)}\ x=2y\\ \textbf{(C)}\ x=60^\circ\\ \textbf{(D)}\ \text{there is no special relationship between }x\text{ and }y\\ \textbf{(E)}\ x=2y\text{ or }x=3y\text{, depending upon the length of }AB$


	\item (\textit{AHSME 1955}) Given a geometric sequence with the first term $\neq 0$ and $r \neq 0$ and an arithmetic sequence with the first term $=0$. 
A third sequence $1,1,2\ldots$ is formed by adding corresponding terms of the two given sequences. 
The sum of the first ten terms of the third sequence is: 


 $\textbf{(A)}\ 978\qquad\textbf{(B)}\ 557\qquad\textbf{(C)}\ 467\qquad\textbf{(D)}\ 1068\\ \textbf{(E)}\ \text{not possible to determine from the information given}$


	\item (\textit{AHSME 1955}) The graphs of $2x+3y-6=0, 4x-3y-6=0, x=2$, and $y=\frac{2}{3}$ intersect in: 


 $\textbf{(A)}\ \text{6 points}\qquad\textbf{(B)}\ \text{1 point}\qquad\textbf{(C)}\ \text{2 points}\qquad\textbf{(D)}\ \text{no points}\\ \textbf{(E)}\ \text{an unlimited number of points}$


	\item (\textit{AHSME 1955}) The expressions $a+bc$ and $(a+b)(a+c)$ are: 


 $\textbf{(A)}\ \text{always equal}\qquad\textbf{(B)}\ \text{never equal}\qquad\textbf{(C)}\ \text{equal whenever }a+b+c=1\\ \textbf{(D)}\ \text{equal when }a+b+c=0\qquad\textbf{(E)}\ \text{equal only when }a=b=c=0$


	\item (\textit{AHSME 1955}) Given $\triangle$ $ABC$ with medians $AE, BF, CD$; $FH$ parallel and equal to $AE$; $BH and HE$ are drawn; $FE$ extended meets $BH$ in $G$. 
Which one of the following statements is not necessarily correct? 


 $\textbf{(A)}\ AEHF\text{ is a parallelogram}\qquad\textbf{(B)}\ HE=HG\\ \textbf{(C)}\ BH=DC\qquad\textbf{(D)}\ FG=\frac{3}{4}AB\qquad\textbf{(E)}\ FG\text{ is a median of triangle }BFH$


	\item (\textit{AHSME 1955}) The graphs of $y=\frac{x^2-4}{x-2}$ and $y=2x$ intersect in: 


 $\textbf{(A)}\ \text{1 point whose abscissa is 2}\qquad\textbf{(B)}\ \text{1 point whose abscissa is 0}\\ \textbf{(C)}\ \text{no points}\qquad\textbf{(D)}\ \text{two distinct points}\qquad\textbf{(E)}\ \text{two identical points}$


	\item (\textit{AHSME 1955}) In order to pass $B$ going $40$ mph on a two-lane highway, $A$, going $50$ mph, must gain $30$ feet. 
Meantime, $C, 210$ feet from $A$, is headed toward him at $50$ mph. If $B$ and $C$ maintain their speeds, 
then, in order to pass safely, $A$ must increase his speed by: 


 $\textbf{(A)}\ \text{30 mph}\qquad\textbf{(B)}\ \text{10 mph}\qquad\textbf{(C)}\ \text{5 mph}\qquad\textbf{(D)}\ \text{15 mph}\qquad\textbf{(E)}\ \text{3 mph}$


\end{enumerate}